\subsection{希尔伯特空间上自同态的交换星代数}

定义 3. --- 设 A 是 \(\mathbf{C}\) 上的代数，E 是复拓扑向量空间。设 π 是 A 在 E 中的一个表示。若对 a ∈ A 的元素 π(a)x 所成的集在 E 中是全的，则称 E 的元素 x 是 π 的循环向量。

设 \(u \in \mathcal{L}(\mathrm{E})\) 是 E 的自同态。我们称 \(x \in \mathrm{E}\) 是 \(u\) 的循环向量，如果它是由 \(u\) 在 \(\mathcal{L}(\mathrm{E})\) 中所生成的星子代数的恒等表示的循环向量。

命题 8. --- 设 E 是复希尔伯特空间。设 A 是 \(\mathcal{L}(\mathrm{E})\) 的交换含幺星子代数。设 \(x\) 是 E 的元素。令 \(\mathrm{E}_x = \overline{\mathrm{A} \cdot x}\) ，并用 \(\mu_x\) 表示 \(x\) 关于 A 的谱测度。

存在唯一的等距同构 \(\theta_x\) ——从 \(\mathrm{L}^2 (\mathsf{X}(\mathrm{A}),\mu_x)\) 到 \(\mathrm{E}_x\) 上——使得 \(\theta_{x}(f) = \mathcal{H}_{\mathrm{A}}(f)(x)\) 对每个函数 \(f\in \mathcal{C}(\mathsf{X}(\mathrm{A}))\) 成立。对每个元素 \(a\in \mathrm{A}\) ，空间 \(\mathrm{E}_x\) 在 \(a\) 下稳定，且若用 \(a_{x}\) 表示由 a 过渡到子空间而得的 \(\mathrm{E}_x\) 的自同态，则我们有 \(a_{x}\circ \theta_{x} = \theta_{x}\circ m_{\mathcal{G}_{\mathrm{A}}(a)}\) 。

设 \(\widetilde{\theta}_x\) 是从 \(\mathcal{C}(\mathrm{X(A)})\) 到 \(\mathrm{E}_x\) 中的线性映射，定义为

\[
\widetilde {\theta} _ {x} (f) = \mathcal {H} _ {\mathrm{A}} (f) (x).
\]

对每个函数 \(f \in \mathcal{C}(\mathsf{X}(\mathsf{A}))\) ，我们有

\[
\| \widetilde {\theta} _ {x} (f) \| ^ {2} = \langle \mathcal {H} _ {\mathrm{A}} (f) x | \mathcal {H} _ {\mathrm{A}} (f) x \rangle = \langle x | \mathcal {H} _ {\mathrm{A}} (| f | ^ {2}) x \rangle = \int_ {\mathsf {X} (\mathrm{A})} | f | ^ {2} d \mu_ {x}.
\]

由于 \(\mathcal{C}(\mathrm{X(A)})\) 在 \(\mathcal{L}^2 (\mathrm{X(A)},\mu_x)\) 中稠密，存在唯一的从 \(\mathrm{L}^2 (\mathrm{X(A)},\mu_x)\) 到 \(\mathrm{E}_x\) 中的等距线性映射，它延拓 \(\widetilde{\theta}_x\) ；我们把它记作 \(\theta_{x}\) 。映射 \(\theta_{x}\) 是满射的，因为它的像是闭的（I, p. 107 的引理 8）且包含 \(\mathrm{A}\cdot x\) 。因而它是一个等距同构。

设 \(a \in \mathrm{A}\) 。令 \(g = \mathcal{G}_{\mathrm{A}}(a) \in \mathcal{C}(\mathsf{X}(\mathrm{A}))\) 。空间 \(\mathrm{E}_x\) 在 \(a\) 下稳定。对 \(f \in \mathcal{C}(\mathsf{X}(\mathrm{A}))\) ，我们有

\[
\widetilde {\theta} _ {x} (m _ {g} (f)) = \mathcal {H} _ {\mathrm{A}} (\mathcal {G} _ {\mathrm{A}} (a) f) x = (a \circ \mathcal {H} _ {\mathrm{A}} (f)) x = a (\widetilde {\theta} _ {x} (f)),
\]

由此得出 \(\theta_x \circ m_g = a_x \circ \theta_x\) 。

推论. — 设 E 是复希尔伯特空间。设 A 是 \(\mathcal{L}(\mathrm{E})\) 的交换含幺星子代数，且具有循环向量 \(x\) 。设 \(\mu_x\) 是 \(x\) 关于 A 的谱测度。存在唯一的等距同构

\[
\theta_ {x} \colon \mathrm{L} ^ {2} (\mathsf {X} (\mathrm{A}), \mu_ {x}) \to \mathrm{E}
\]

使得 \(\theta_x(f) = \mathcal{H}_{\mathrm{A}}(f)\) 对每个函数 \(f \in \mathcal{C}(\mathsf{X}(\mathrm{A}))\) 成立。对 A 中每个 a，我们有 \(a \circ \theta_x = \theta_x \circ m_{\mathcal{G}_{\mathrm{A}}(a)}\) 。

事实上，这时我们有 \(E_{x}=E\) ，且 \(a_{x}=a\) （对所有 \(a\in A\) ）。

设 E 是复希尔伯特空间。设 \(x \in E\) ，A 是 \(\mathcal{L}(\mathrm{E})\) 的交换含幺星子代数。用 \(E_{x}\) 表示 E 的闭子空间 \(\overline{A \cdot x}\) 。对每个 \(a \in A\) ，我们有 \(a(\mathrm{E}_{x}) \subset \mathrm{E}_{x}\) 。用 \(A_{x}\) 表示 \(\mathcal{L}(\mathrm{E}_{x})\) 的交换含幺星子代数，它由 A 的元素经过渡到子空间而得的 \(E_{x}\) 的自同态组成。向量 x 是 \(A_{x} \subset \mathcal{L}(\mathrm{E}_{x})\) 的循环向量。

命题 9. --- 存在 E 的子集 C，使得 E 是诸空间 \(E_{x}\) （\(x \in C\)）的希尔伯特直和。若 E 是可数生成的，则集 C 是可数的。

设 \(\mathcal{O}\) 是由 E 的这样的子集 C 组成的集：诸子空间 \(\mathrm{E}_x\) （\(x\in \mathrm{C}\)）两两正交。集 \(\mathcal{O}\) 按包含关系有序，它是有限特征的（E, III, p. 34，定义 2），因为 C 属于 \(\mathcal{O}\) 当且仅当由 C 的两个元素组成的集属于 \(\mathcal{O}\) 。由 E, III, p. 35 的定理 1，存在 \(\mathcal{O}\) 的一个极大元 C。

设 F 是由诸子空间 \(E_{x}\) （\(x \in C\)）生成的 E 的闭子空间。只须证明 \(F^{\circ}\) 归结为 0 即可完成命题的证明。设 y 是 \(F^{\circ}\) 的元素。对每个 \(x \in C\) 及 A 中所有 a 与 b，我们有 \(\langle a(y) | b(x) \rangle = \langle y | a^{*}b(x) \rangle = 0\) ，因为 \(a^{*}b(x)\) 属于 \(E_{x}\) ，从而属于 F。由于元素 \(a(y)\) （相应地， \(b(x)\) ）生成 \(E_{y}\) （相应地， \(E_{x}\) ）的稠密子空间，我们有 \(E_{y} \subset E_{x}^{\circ}\) 。由于 C 在 O 中是极大的，这就意味着 \(E_{y}\) 归结为 0，从而 y = 0。

假设 E 是可数生成的。由于 \(E_{x}\) 对所有非零 \(x \in C\) 非零，任何使 E 成为诸空间 \(E_{x}\) （\(x \in C\)）的希尔伯特直和的集 C 都是可数的。

定理 1. --- 设 A 是 \(\mathcal{L}(\mathrm{E})\) 的交换含幺星子代数。存在局部紧拓扑空间 X、X 上的正测度 \(\mu\) 、等距同构 \(\theta\) ——从 \(\mathrm{L}^2 (\mathrm{X},\mu)\) 到 E 上——以及星代数等距态射 \(\pi\) ——从 A 到 \(\mathcal{C}_b(\mathrm{X})\) 中——使得对每个 \(a\in \mathrm{A}\) ，有

\[
a \circ \theta = \theta \circ m _ {\pi (a)}.
\]

由命题 9，存在 E 的子集 C，使得 E 是诸子空间 \(E_{x}\) （\(x \in C\)）的希尔伯特直和。设 X 为局部紧拓扑空间 \(\mathsf{X}(\mathsf{A}) \times \mathsf{C}\) ，其中 C 赋有离散拓扑。存在 X 上唯一的测度 \(\mu\) ，使得 \(\mu|(\mathsf{X}(\mathsf{A}) \times \{x\})\) 与 x 的谱测度 \(\mu_{x}\) 等同（对每个 \(x \in C\) ；INT, III, p. 65, § 2, n°1，命题 1）；这个测度是正的（参见同前）。于是有希尔伯特直和分解

\[
\mathrm{L} ^ {2} (\mathrm{X}, \mu) = \bigoplus_ {x \in \mathrm{C}} \mathrm{L} ^ {2} (\mathrm{X} (\mathrm{A}), \mu_ {x})
\]

（IV, p. 179 的命题 1, c) 应用于诸集 \(\mathsf{X}(\mathsf{A}) \times \{x\}\) （\(x \in \mathrm{C}\)））。因此存在唯一的等距 \(\theta: \mathrm{L}^2(\mathrm{X}, \mu) \to \mathrm{E}\) ，它与 \(\theta_x: \mathrm{L}^2(\mathsf{X}(\mathsf{A}), \mu_x) \to \mathrm{E}_x\) 在 \(\mathrm{L}^2(\mathsf{X}(\mathsf{A}), \mu_x)\) 上重合（对所有 \(x \in \mathrm{C}\) ）。

设 \(p \colon \mathrm{X} \to \mathrm{X}(\mathrm{A})\) 是典范投影；它是满射的，因而线性映射 \(p^* \colon \mathcal{C}_b(\mathrm{X}(\mathrm{A})) \to \mathcal{C}(\mathrm{X})\) ——由 \(f \mapsto f \circ p\) 所定义——是单射的。令 \(\pi = p^* \circ \mathcal{G}_{\mathrm{A}}\) ；这是从 A 到 \(\mathcal{C}_b(\mathrm{X})\) 中的星代数的单射态射；因而它是等距的（I, p. 112 的命题 9）。

设 \(a \in \mathrm{A}\) 。对每个 \(x \in \mathrm{C}\) ，我们有 \(a_x \circ \theta_x = \theta_x \circ m_{\mathcal{G}_A(a)}\) 。于是连续线性映射 \(a \circ \theta\) 与 \(\theta \circ m_{\pi(a)}\) 在诸子空间 \(\mathrm{L}^2(\mathrm{X}(\mathrm{A}), \mu_x)\) 上重合，因而相等。

推论（谱定理）. --- 设 E 是复希尔伯特空间。设 \(u \in \mathcal{L}(\mathrm{E})\) 是 E 的正规自同态。存在局部紧拓扑空间 X、X 上的正测度 \(\mu\) 、等距同构 \(\theta\) ——从 \(\mathrm{L}^{2}(\mathrm{X},\mu)\) 到 E 上——以及 X 上的连续有界函数 g，使得 \(u = \theta \circ m_{g} \circ \theta^{-1}\) 。

若 u 有循环向量 x，则可取 \(X = \mathrm{Sp}(u)\) ，并取 g 为 X 的恒等函数。

只须对 \(\mathcal{L}(\mathrm{E})\) 的、由 \(u\) 所生成的星子代数 A 应用定理 1（相应地，命题 8 的推论），并令 \(g = \pi(u)\) 。

这个命题把关于希尔伯特空间单个正规自同态的每个问题化归为关于乘法算子的类似问题，这往往使其研究得到简化（例如参见 IV, p. 325 的习题 19）。

